\subsection{Training an Encoder from Scratch on Raw Pixels (No Pretrained Features)}
\label{exp:raw-pixel-from-scratch}

\paragraph{What we tried} Every checkpoint logged so far starts from frozen or
fine-tuned DINOv2 features. \texttt{pwm\_raw\_residual\_mlp\_paper\_window0\_vicreg}
removes that dependency entirely: \texttt{token\_mode="raw"} trains a small
residual-MLP encoder directly on raw pixels, from scratch, jointly with the predictor
--- otherwise the same \texttt{paper\_window0\_vicreg} recipe (VICReg regularizer,
\texttt{window\_size=0}, \texttt{straighten\_weight=0.1}, \texttt{batch\_size=32}) used
elsewhere in this project. This tests whether the patterns found across this
investigation depend on starting from a strong pretrained visual prior, or persist even
when the encoder has to build its own representation from nothing. Scored with the same
full metric suite from Section~\ref{sec:methodology} (Main and Additional).

\paragraph{Result}

\begin{table}[H]
\centering
\caption{Main metrics, raw-pixel residual-MLP encoder (from scratch), K400.}
\label{tab:raw-pixel-main}
\begin{tabular}{lc}
\toprule
Metric & Value \\
\midrule
\multicolumn{2}{l}{\textit{Whole-dataset collapse and spread}} \\
effective\_rank\_fraction & \good{0.8419} \\
$W$ (within-clip mean pairwise L2) & 14.9135 \\
$B$ (across-clip mean pairwise L2) & 26.7479 \\
$B-W$ (gap) & 11.8344 \\
\midrule
\multicolumn{2}{l}{\textit{Trajectory geometry}} \\
trajectory\_straightness ($z$) & \good{0.2090} \\
mean\_velocity\_cosine ($z$) & \good{0.0841} \\
start\_end\_displacement ($z$) & 0.9174 \\
direction\_reversal\_fraction ($x$) & \good{0.4083} \\
\midrule
\multicolumn{2}{l}{\textit{Temporal order}} \\
time\_progress\_spearman ($z$) & 0.9070 \\
distance\_to\_goal\_spearman ($x$) & \good{0.6425} \\
adjacent\_temporal\_nn\_fraction ($z$) & 0.8699 \\
\midrule
\multicolumn{2}{l}{\textit{Cross-clip temporal consistency}} \\
tcc\_kendall\_tau (mean) & 0.0136 \\
\midrule
\multicolumn{2}{l}{\textit{Semantic information}} \\
chirality.mean\_accuracy ($x$) & \bad{0.5811} \\
content.knn.balanced\_acc & 0.0623 \\
content.linear.balanced\_acc & 0.0772 \\
change.knn.balanced\_acc & 0.1057 \\
change.linear.balanced\_acc & 0.1180 \\
content+change.knn.balanced\_acc & 0.1159 \\
content+change.linear.balanced\_acc & 0.1354 \\
\bottomrule
\end{tabular}

\vspace{0.3em}
{\footnotesize \texttt{content.linear.balanced\_acc} (0.0772) is below \texttt{chance.majority}
(0.1280) --- the only checkpoint in this project's log where the linear content probe
underperforms the majority-class baseline. No \texttt{reconstruction\_r2} row: there is
no DINOv2 target to reconstruct back to for a \texttt{token\_mode="raw"} checkpoint (see
\texttt{eval\_common/loading.py}'s \texttt{encode\_world\_model} docstring).}
\end{table}

\begin{table}[H]
\centering
\caption{Additional metrics, raw-pixel residual-MLP encoder --- same run as
Table~\ref{tab:raw-pixel-main}.}
\label{tab:raw-pixel-additional}
\begin{tabular}{lc}
\toprule
Metric & Value \\
\midrule
\multicolumn{2}{l}{\textit{Trajectory geometry and raw variants}} \\
mean\_chord\_alignment ($z$) & 0.1530 \\
trajectory\_straightness\_raw ($x$) & 0.2089 \\
mean\_velocity\_cosine\_raw ($x$) & 0.0841 \\
mean\_chord\_alignment\_raw ($x$) & 0.1529 \\
\midrule
\multicolumn{2}{l}{\textit{Raw embedding magnitude}} \\
norm\_growth ($x$) & 1.0003 \\
norm\_std\_within\_clip ($x$) & 0.1088 \\
adjacent\_l2\_distance ($x$) & 6.9463 \\
mean\_norm ($x$) & 20.8358 \\
\midrule
\multicolumn{2}{l}{\textit{Change-direction consistency}} \\
change\_direction.gap & 0.0088 \\
\bottomrule
\end{tabular}

\vspace{0.3em}
{\footnotesize The normalized ($z$) and raw ($x$) trajectory-geometry values are nearly
identical here (e.g.\ \texttt{trajectory\_straightness} 0.2090 vs.\ 0.2089) --- this
encoder's per-frame norm is already close to constant within a clip
(\texttt{norm\_std\_within\_clip}=0.1088, \texttt{norm\_growth}$\approx$1.0003), so
L2-normalizing barely changes the geometry.}
\end{table}

\begin{table}[H]
\centering
\caption{Training stats, raw-pixel residual-MLP encoder.}
\label{tab:raw-pixel-training}
\begin{tabular}{lccc}
\toprule
Loss term & Epoch 0 & Best (ep17) & Last (ep49) \\
\midrule
train\_pred & 0.316 & 0.212 ($\downarrow$33\%) & 0.327 \\
train\_straighten & 1.387 & 1.124 ($\downarrow$19\%) & 1.110 \\
train\_vicreg\_std & 0.307 & 0.238 ($\downarrow$22\%) & 0.135 \\
train\_vicreg\_cov & 2.812 & 3.547 ($\uparrow$26\%, worse) & 3.811 \\
train\_total & 0.874 & 0.704 ($\downarrow$20\%) & 0.726 \\
\bottomrule
\end{tabular}

\vspace{0.3em}
{\footnotesize Best epoch 17 of 50 logged; full run, no early stop.}
\end{table}

\paragraph{Inference}

\textbf{This pushes the generalization question from the previous experiment
further --- not just a different encoder \emph{architecture}, but no pretrained visual
prior at all.} The batch-size/regularizer ablation showed the collapse-vs-straightness
tension isn't specific to the patch-channel encoder; this checkpoint asks whether it
depends on starting from frozen DINOv2 features in the first place, by training a small
residual-MLP encoder from scratch directly on raw pixels, jointly with the predictor,
using the same \texttt{paper\_window0\_vicreg} recipe used elsewhere in this log.

\textbf{This is the only checkpoint in the log so far to be green on collapse
\emph{and} every straightness metric simultaneously --- everywhere else in this log,
those two have traded off.} \texttt{effective\_rank\_fraction}=0.8419 beats both DINOv2
baselines (green), and \texttt{trajectory\_straightness} (0.209), \texttt{mean\_velocity\_cosine}
(+0.0841, the only \emph{positive} value seen anywhere in this log --- every other
checkpoint and every frozen encoder is negative), and \texttt{direction\_reversal\_fraction}
(0.408, the lowest/best of anything logged) are all green too. In every prior experiment,
a green collapse metric came with red straightness metrics or vice versa; here all four
are green at once. \texttt{distance\_to\_goal\_spearman} (0.6425) is green too, beating
both DINOv2 baselines --- \texttt{time\_progress\_spearman} (0.9070) is the only one of
the eight core metrics that lands uncolored, close to but just under both baselines
(0.9040--0.9114). This is the best-rounded profile of any trained checkpoint in the log:
not just collapse and straightness together, but one of the two temporal-order metrics
as well.

\textbf{That combination does not extend to semantic content --- \texttt{content.linear.balanced\_acc}
is the one clearly red result, and it undercuts even chance.} 0.0772, below this
checkpoint's own \texttt{chance.majority} baseline of 0.1280 --- the only checkpoint in
this project's log where the linear content probe underperforms just guessing the
majority class. \texttt{chirality.mean\_accuracy} (0.5811) is also red, in the same
range as every other trained checkpoint in this log, not recovered by the otherwise
exceptional collapse/straightness profile. With no pretrained visual prior constraining
it, the encoder is free to organize its space however is easiest for the loss --- and
the easiest solution found here is a very straight, very ordinal, non-collapsed
trajectory that carries almost no recoverable content about what the clip actually
shows. This is the clearest single data point in the log for why collapse and
straightness numbers, even both green at once, cannot be read as "this representation is
good" without also checking a semantic probe.

\textbf{Training dynamics are a mild drift, not the sharp divergence seen when
\texttt{window\_size=0} destabilized the patch-channel encoder.} \texttt{train\_total}
bottoms out smoothly at epoch 17 (0.704) and creeps back up gradually to 0.726 by epoch
49 --- never exploding the way \texttt{window0\_2}/\texttt{window0\_vicreg\_3} did in the
window/VICReg ablation. \texttt{train\_vicreg\_cov} never improves at any point in
training, including the "good" early phase (2.81 $\to$ 3.55 $\to$ 3.81, monotonically
worse throughout) --- the early gains (epoch 0 $\to$ 17) come entirely from
\texttt{train\_vicreg\_std} and \texttt{train\_straighten}, not from \texttt{vicreg\_cov}
doing any real work here, yet \texttt{effective\_rank\_fraction} still ends up green ---
a reminder that a single loss term's own trajectory doesn't necessarily predict the
eval-time collapse metric it's meant to serve.
